\documentclass[12pt]{article}
\usepackage[utf8]{inputenc}
\usepackage[T1]{fontenc}
\usepackage{hyperref}
\title{Kamba (Brazil) Language Resources for K-12 Curriculum Development}
\author{Compiled by Claude Sonnet 5.5 (claude-sonnet-5-5) for Andr\'as Kornai}
\date{}
\begin{document}
\maketitle

\section{Introduction}
The code xba (``Kamba (Brazil)'', Glottocode kamb1287) does not appear to denote a distinct language. According to the Glottolog entry, xba was listed as an extinct, unclassified language of Mato Grosso do Sul that ``may have been Tup\'i''. Glottolog states that the Kamba identify as Chiquitano and that the attested language is clearly Chiquitano (ISO cax). The code was reportedly retired as a duplicate of Chiquitano in 2016. The family label ``Bookkeeping'' in the request is consistent with a bookkeeping entry rather than a real language family. A native-language K-12 curriculum for ``Kamba (Brazil)'' therefore cannot be built from xba-specific materials. Chiquitano resources would be the relevant starting point, and I did not search for them.

\section{Primary Educational Resources}
I found no primers, textbooks or school programmes specific to xba in the search results.

\section{Linguistic and Reference Materials}
The Glottolog languoid page for Kamba (Brazil) documents the retirement and the Chiquitano identification: \url{https://glottolog.org/resource/languoid/id/kamb1287}.

\section{Religious and Community Resources}
Joshua Project lists a Kamba people in Brazil of about 2,300 people, noting that their presence is not officially recognised: \url{https://joshuaproject.net/people_groups/19152/BR}. Wikipedia has a general article on the Camba (people of Bolivian origin, in the Corumb\'a area of Mato Grosso do Sul according to the scholarly work I saw): \url{https://en.wikipedia.org/wiki/Camba}. I found no scripture translations or community language organisations for xba.

\section{Digital and Audio Resources}
None found for xba.

\section{Curriculum Development Context}
There is no evidence of a separate xba language to teach. Any curriculum effort for this community should first clarify with the community whether it concerns Chiquitano, Spanish (the language the community is reported to speak), or Portuguese, and should consult Chiquitano-specific resources.
\end{document}